\chapter{提案手法}
% 【この章の目的】第三者が再実装できる粒度で，手法の全体像と各要素を記述する．
% 内容入り例：examples/chapters/04-method-example.tex
% まず処理フロー全体を示し，その後に各構成要素を説明する．

\section{全体構成}
% 入出力，主要モジュール，学習時と推論時の違いを図示するとよい．図題は図の下へ置き，本文中で参照する．

\section{主要な処理}
% 数式と文章を対応させ，各変数の意味，設計理由，計算手順を示す．番号付き数式は式\eqref{eq:...}として本文から参照する．

\section{実装条件}
% 前処理，ハイパーパラメータ，停止条件，使用ライブラリ等を記載する．
